\documentclass[12pt,a4paper]{article}
\usepackage[utf8]{inputenc}
\usepackage{geometry}
\geometry{a4paper, margin=1in}
\usepackage{booktabs}
\usepackage{graphicx}
\usepackage{hyperref}
\usepackage{amsmath}

\title{Multi-GPU Parallelism Strategies for Micromagnetic Minimizers in JAX}
\author{}
\date{}

\begin{document}
\maketitle

\section{Introduction}
Based on the \texttt{tr} minimizer and \texttt{poisson\_solve.py} source code, we can make a highly accurate, realistic estimate of communication overhead versus compute speedup. The fundamental difference between parallelization strategies lies in \textbf{where the communication happens relative to the algorithmic loops}.

\section{Assumptions for a Realistic Large Mesh}
\begin{itemize}
    \item \textbf{Nodes ($N$):} 10 Million ($10^7$) nodes.
    \item \textbf{Data Type:} \texttt{float64} (8 bytes per scalar).
    \item \textbf{Scalar vector size} (e.g., $U$ or $p$): $N \times 8 \text{ bytes} \approx \textbf{80 MB}$.
    \item \textbf{Vector field size} (e.g., $m$ or $v$): $3N \times 8 \text{ bytes} \approx \textbf{240 MB}$.
    \item \textbf{Inner TR Iterations ($K_{TR}$):} 5 iterations.
    \item \textbf{Poisson PCG Iterations ($K_{CG}$):} 100 iterations (a conservative estimate for a preconditioned Poisson solve).
    \item \textbf{Compute Baseline:} $\sim$2.5 ms for one full matrix-vector product (SpMV) on a single GPU.
    \item \textbf{Line Search:} \texttt{pcohen\_hs} averages \textbf{2 trial steps} (evaluating Poisson twice). \texttt{tr} evaluates Poisson \textbf{exactly 1 time}.
    \item \textbf{Hardware:} 4 GPUs. PCIe Gen4 ($\sim$25 GB/s) or NVLink ($\sim$300 GB/s).
\end{itemize}

\section{Optimal GPU Assignment for Method 1 (Task Parallelism)}
The beauty of Method 1 is that the optimal mapping of sparse matrices to GPUs is \textbf{universal}. It remains exactly the same regardless of whether you use PCIe or NVLink, and regardless of whether you use the \texttt{tr} or \texttt{pcohen\_hs} minimizer. The goal is to balance the memory load and isolate the heavy inner loops so they do not compete for hardware resources.

To maximize concurrency, you must use the split spatial components (\texttt{Dx\_sparse, Dy\_sparse, Dz\_sparse} and \texttt{Gx\_sparse, Gy\_sparse, Gz\_sparse}) rather than the concatenated \texttt{D} and \texttt{G} matrices.

\subsection*{Recommended 4-GPU Mapping}
\begin{itemize}
    \item \textbf{GPU 0:} \texttt{A\_sparse} (Poisson Stiffness)
    \item \textbf{GPU 1:} \texttt{K\_eff\_sparse} (Exchange/Anisotropy), \texttt{Dx\_sparse}, \texttt{Gx\_sparse}
    \item \textbf{GPU 2:} \texttt{Dy\_sparse}, \texttt{Gy\_sparse}
    \item \textbf{GPU 3:} \texttt{Dz\_sparse}, \texttt{Gz\_sparse}
\end{itemize}

\subsection*{Why this mapping is optimal for all combinations:}
\begin{enumerate}
    \item \textbf{Hardware Isolation:} The Poisson PCG loop (\texttt{A\_sparse}) is the most memory-bandwidth intensive operation. By dedicating GPU 0 solely to \texttt{A\_sparse}, the PCG loop gets 100\% of GPU 0's resources without being interrupted.
    \item \textbf{Hidden Physics (Inner Iterations):} Both \texttt{tr} and \texttt{pcohen\_hs} rely on an inner preconditioner/Newton loop that exclusively uses \texttt{K\_eff\_sparse}. By placing \texttt{K\_eff\_sparse} on GPU 1, it runs on completely separate hardware from GPU 0. Furthermore, when evaluating the final energy, JAX will schedule \texttt{K\_eff\_sparse} (on GPU 1) to run \textit{simultaneously} in the background while the Poisson PCG solver runs on GPU 0.
    \item \textbf{Perfect Spatial Concurrency:} During the setup phase (computing the RHS), GPUs 1, 2, and 3 will fire simultaneously to compute \texttt{Dx}, \texttt{Dy}, and \texttt{Dz}. After the Poisson solve finishes, they will fire simultaneously again to compute the Demag gradients \texttt{Gx}, \texttt{Gy}, and \texttt{Gz}.
    \item \textbf{Memory Balance:} \texttt{K\_eff\_sparse} is the largest matrix ($3N \times 3N$). GPUs 1, 2, and 3 only hold tiny $N \times N$ directional matrices, meaning GPU 1 has plenty of free VRAM to comfortably host \texttt{K\_eff\_sparse}.
\end{enumerate}

\section{Performance Estimates per Outer Minimizer Step}

\begin{table}[h]
\centering
\resizebox{\textwidth}{!}{%
\begin{tabular}{lllrrlrl}
\toprule
\textbf{Minimizer} & \textbf{Method} & \textbf{Interconnect} & \textbf{Comm. Vol.} & \textbf{Comm. Time} & \textbf{Compute} & \textbf{Total Time} & \textbf{Summary \& Bottlenecks} \\ \midrule
\texttt{tr} & Task (M1) & PCIe & 0.63 GB & 26 ms & 270 ms & \textbf{$\sim$296 ms} & \textbf{Best for PCIe.} Sequential compute, low comms. \\
\texttt{tr} & Task (M1) & NVLink & 0.63 GB & 2 ms & 270 ms & \textbf{$\sim$272 ms} & Fails to use NVLink to speed up computation. \\
\texttt{tr} & Data (M2) & PCIe & 9.53 GB & 390 ms & 68 ms & \textbf{$\sim$458 ms} & \textbf{AVOID.} 4x speedup destroyed by 10GB PCIe transfer. \\
\texttt{tr} & Data (M2) & NVLink & 9.53 GB & 33 ms & 68 ms & \textbf{$\sim$101 ms} & \textbf{Best for NVLink.} Massive 4x compute speedup. \\ \midrule
\texttt{pcohen} & Task (M1) & PCIe & 1.25 GB & 51 ms & 528 ms & \textbf{$\sim$579 ms} & Takes $\sim$2x longer due to line search Poisson solves. \\
\texttt{pcohen} & Task (M1) & NVLink & 1.25 GB & 4 ms & 528 ms & \textbf{$\sim$532 ms} & Bound by sequential compute of the line search. \\
\texttt{pcohen} & Data (M2) & PCIe & 17.89 GB & 733 ms & 132 ms & \textbf{$\sim$865 ms} & \textbf{WORST.} Line search forces $\sim$18 GB transfer. \\
\texttt{pcohen} & Data (M2) & NVLink & 17.89 GB & 61 ms & 132 ms & \textbf{$\sim$193 ms} & Fast, but $\sim$2x slower than \texttt{tr} on NVLink. \\ \bottomrule
\end{tabular}%
}
\caption{Estimated performance metrics across different configurations.}
\end{table}

\section{Key Takeaways}

\begin{enumerate}
    \item \textbf{Hardware dictates the Method:}
    \begin{itemize}
        \item \textbf{No NVLink (PCIe only):} You are strictly bound to \textbf{Method 1 (Task Parallelism / 1 Matrix per GPU)}. Any attempt to use \texttt{all\_gather} across the rows (Method 2) will saturate the motherboard and slow your code down dramatically.
        \item \textbf{With NVLink:} You must use \textbf{Method 2 (Data Parallelism)}. Method 1 will leave 75\% of your immense compute power sitting idle.
    \end{itemize}

    \item \textbf{Algorithm dictates the Baseline Time:}
    \begin{itemize}
        \item \textbf{\texttt{tr} minimizes Poisson solves:} Because Trust-Region relies on its inner loop (which uses only $K_{eff}$ and is completely isolated from the Poisson equation) and evaluates the actual energy exactly \textit{once}, it is inherently faster per-step than \texttt{pcohen\_hs}.
        \item \textbf{\texttt{pcohen\_hs} doubles communication and compute:} The Armijo line search is a brute-force exploratory loop. If it takes 2 trial steps to find a descent direction, it has to completely solve the Poisson equation twice, doubling the PCG iterations and doubling the massive cross-GPU data transfers.
    \end{itemize}
\end{enumerate}

\section{Final Verdict}
If you want the fastest possible simulation that fits in memory on a multi-GPU node \textit{without} NVLink, use \textbf{\texttt{tr} + Method 1} (and implement the \texttt{energy\_only} scalar reduction trick to lower the 26 ms communication time even further). If you have NVLink, use \textbf{\texttt{tr} + Method 2}.

\end{document}
